% Candidate A.60. Notes only until independent audits and checkpoint release.
\subsection{The next Heisenberg cover: an exact seven-adic model}

\begin{statement}{A.60}{Proved}{Both all-prime assertions for the genus-one-hundred-forty-eight Heisenberg cover}
Let \(E:y^2=x^3-x\), and let \(\psi _7\) be its seventh
division polynomial. Let \(C_7\) denote the smooth proper
normalization of the displayed function field. The actual
curve and punctured map
\[
 C_7:\ y^2=x^3-x,\ z^7=\psi _7(x),\qquad
 Y_7=C_7\setminus q^{-1}(E[7])
              \xrightarrow{[7]q}E\setminus\{O\}       \tag{A.60.1}
\]
give a geometrically connected Galois cover with Heisenberg
group of order \(343\), proper genus \(148\), and \(49\)
cusps. For every coefficient prime, its nontrivial closed
procyclic subgroups with open arithmetic normalizer are
cuspidal. For every selected cusp packet and every indicated
open derived subgroup, its full geometric pro-prime fixed
subgroup is exactly the selected cusp inertia.

The fixed-group assertion comes from an explicit seven-adic
compact-type model. Its positive-genus components are one
supersingular elliptic main component and seven copies of
\[
                              V^7-V=U^8,               \tag{A.60.2}
\]
each of genus twenty-one. All forty-nine cusps lie on a
rational marked tree attached to the main component, away
from the seven positive-genus tails.
\end{statement}

\begin{proof}
We first specify the actual cover and its group, so the
model calculation concerns an identified instance of Task A.5.
Use the exact polynomials
\[
\begin{aligned}
 \psi _3&=3x^4-6x^2-1,\\
 B&=x^6-5x^4-5x^2+1,\\
 \psi _7&=\psi _5\psi _3^3-128(x^3-x)^2B^3,
 \qquad h=\frac{\psi _7+1}{7},                         \tag{A.60.3}
\end{aligned}
\]
where \(\psi _5\) is the polynomial of A.54. These are the
standard division polynomials of \cite[Exercise III.3.7,
pp.105--106]{fghSilvermanAEC}. The middle
identity is the division-polynomial recurrence
\(\psi _7=\psi _5\psi _3^3-\psi _2\psi _4^3\), with
\(\psi _2=2y\) and \(\psi _4=4yB\).
It gives \(h\in\Z[x]\), and
\[
 \overline h=x^{24}-2x^{22}-2x^{20}+x^{18}-x^{14}
                  -3x^{10}-x^8-x^6-3x^4\quad\text{in }\F_7[x].
                                                        \tag{A.60.4}
\]

\paragraph{The actual geometric Galois group.}
The divisor is
\(\operatorname{div}(\psi _7)=D_{E[7]}-49(O)\),
where \(D_{E[7]}\) includes \(O\); see
\cite[Exercise III.3.7(f), p.106]{fghSilvermanAEC}. Its order at each
nonzero seven-torsion point is one, and its order at
\(O\) is \(-48\), also prime to seven. Thus the
Kummer cover \(q\) is connected of degree seven and
totally ramified precisely at the forty-nine points of
\(E[7]\). Hurwitz gives
\[
                2g(C_7)-2=49(7-1)=294,
                           \qquad g(C_7)=148.          \tag{A.60.5}
\]
The characteristic-zero isogeny \([7]\) is etale, so
(A.60.1) has degree \(343\), has forty-nine cusps and
is etale away from them.

Every translation by \(T\in E[7]\) lifts to \(C_7\)
after finite extension. Indeed the divisor of
\(\psi _7(P+T)/\psi _7(P)\) is
\(49(O)-49(-T)\), seven times a principal divisor
\(7(O)-7(-T)\). Choose a function \(r_T\) for that
divisor and a seventh root of the remaining constant.
Then \((P,z)\mapsto(P+T,a_T r_T(P)z)\) is an actual
lift. These give all \(343\) automorphisms of the
degree-\(343\) composite, with exact sequence
\[
                 1\longrightarrow\mu _7
                 \longrightarrow G\longrightarrow E[7]
                 \longrightarrow1.                   \tag{A.60.6}
\]
The extension is central. It is nonabelian: the original
once-punctured elliptic inertia is a commutator, whose
image in this cover has order seven, whereas it would
vanish in an abelian quotient.

For precision, the group has exponent seven, not
forty-nine. The evenness of \(\psi _7\) gives the actual
automorphism \(\rho(x,y,z)=(-x,iy,z)\), which acts by
\([i]\) on \(E[7]\) and trivially on \(\mu _7\).
In a central extension (A.60.6), seventh powers define
an additive map \(E[7]\to\mu _7\): the commutator
correction has exponent \(\binom72\), divisible by
seven. This map is invariant under \([i]\). Since
\([i]-1\) is invertible on \(E[7]\), it vanishes.
Choose a basis of \(E[7]\) with nonzero commutator
pairing and lift it; together with the central generator
its relations are precisely those of the Heisenberg
group of order \(343\). This also puts the normalizer
assertion within A.59. All these maps and groups descend
after a finite extension, which is sufficient for the
arithmetic normalizer comparison.

\paragraph{The radicial main component and its seven critical centres.}
Choose \(\pi^7=7\) and substitute \(z=-1+\pi u\).
The exact equation is
\[
 u^7-\pi^6u^6+3\pi^5u^5-5\pi^4u^4
             +5\pi^3u^3-3\pi^2u^2+\pi u=h(x),
                \qquad y^2=x^3-x.                     \tag{A.60.7}
\]
Modulo seven it has function field \(k(E)(h^{1/7})\).
The exact differential factorization in characteristic seven is
\[
                   dh=3x^3(x^2-1)^3(x^2-3)^7\,dx\ne0.
                                                        \tag{A.60.8}
\]
Thus this field is \(k(E)^{1/7}\). Its normalized curve
is an elliptic Frobenius twist, and (A.60.7) is smooth
at its special generic point, giving an actual
genus-one main component mapping radicially with degree
seven to \(E\). The elliptic special fibre is
supersingular: the coefficient of \(x^6\) in
\((x^3-x)^3\) is zero.

The zeros of (A.60.8) on \(E\) are exactly
\[
                (0,0),(1,0),(-1,0),
                            \qquad x^2=3,              \tag{A.60.9}
\]
the latter condition giving four points. Every zero
has order seven. At a two-torsion point use \(y\)
as parameter: \(x-a\) has order two and \(dx\)
order one. At each of the four other points, \(x-a\)
is a parameter. These are seven distinct finite points;
none is the origin.

\paragraph{The degree-seven coefficient which must also vanish.}
The differential alone would not exclude a seventh-power
term below the desired degree-eight term. We verify that
term explicitly. Choose the centres (A.60.9) in a finite
\emph{unramified} extension \(R_0/\Z_7\), taking
\(a^2=3\) and a unit root \(y_0^2=a^3-a\) for the
four non-two-torsion points. Formula (A.60.4) gives
\[
 \begin{aligned}
  \overline h(a+t)-\overline h(a)&\equiv2t^8\pmod {t^9},
                                      &&a^2=3,\\
  \overline h(a+t)-\overline h(a)&\equiv b_at^4\pmod {t^5},
                                      &&a=0,1,-1,\\
                         b_0&=4,\qquad b_1=b_{-1}=2.
 \end{aligned}                                          \tag{A.60.10}
\]
In particular every coefficient of degree one through
seven in the first line vanishes, including degree
seven. These are identities in
\(\F_7[a,t]/(a^2-3)\), not only checks at one residue
embedding. The certificate expands the entire
polynomial and verifies these identities exactly.
All corresponding lower coefficients over \(R_0\)
are therefore divisible by \(7\).

\paragraph{Exact smooth algebraic tail charts.}
Choose \(\eta^6=-7\), \(\rho^8=\eta\), and
\(z_0^7=\psi _7(a)\). The latter is a unit congruent
to \(-1\). Use \(z=z_0(1+\eta V)\). The exact
covering equation, after division by \(7\eta\), is
\[
 \begin{split}
 z_0^7\bigl(V+3\eta V^2+5\eta^2V^3+5\eta^3V^4
                         +3\eta^4V^5+\eta^5V^6-V^7\bigr)
              =\frac{h(x)-h(a)}{\eta}.
 \end{split}                                            \tag{A.60.11}
\]
For \(y_0\ne0\), set \(x=a+\rho U\),
\(y=y_0+\rho W\). The elliptic equation is
\[
 2y_0W+\rho W^2
              =(3a^2-1)U+3a\rho U^2+\rho^2U^3.       \tag{A.60.12}
\]
By (A.60.10), the right side of (A.60.11) is
integral and reduces to \(2U^8\). A lower coefficient
of degree \(j<8\) has valuation at least
\(1+(j-8)/48>0\); a higher coefficient has valuation
at least \((j-8)/48>0\). The elliptic equation
eliminates \(W\) by its unit linear coefficient.

At \(a=0,1,-1\), instead set \(y=\rho U\),
\(x=a+\rho^2W\). Put \(\delta=3a^2-1\), a unit.
The first equation becomes
\[
                U^2=\delta W+3a\rho^2W^2+\rho^4W^3.
                                                        \tag{A.60.13}
\]
The second line of (A.60.10) makes the right side
of (A.60.11) integral, with reduction
\(b_aW^4=b_a\overline\delta^{-4}U^8\).
Again all other coefficients have strictly positive
valuation at the indicated scale.

The special equations in both cases are therefore
\(V^7-V=c_aU^8\), with \(c_a\ne0\). The two
equation partials in \(W,V\) are units on the
special fibre, so these are actual smooth algebraic
charts with integral special curve. Their generic
coordinate changes are invertible on \(C_7\).
The Artin--Schreier polynomial is irreducible because
its pole order eight is prime to seven. Over the
algebraic residue field rescaling \(U\) gives
(A.60.2).

For its ramification calculation put \(t=U^{-1}\).
At the unique point above infinity, \(v(t)=7\),
\(v(V)=-8\), and \(t^7V^6\) is a uniformizer.
Under a nontrivial translation \(V\mapsto V+b\),
its difference has valuation \(49-40=9\). The
lower jump is eight and the different is \(6\cdot9\).
Hurwitz gives genus twenty-one. The one-branch-point
Deuring--Shafarevich formula gives seven-rank zero,
as in \cite[Corollary 1.8]{fghCrew1984}.

\paragraph{Stable components and all forty-nine marked sections.}
The actual positive-genus valuations furnished by
these seven smooth charts persist in the stable
model, by the regular-surface factorization and
positive-genus persistence argument of A.52 and
A.58; the precise sources are
\cite[Theorem 9.2.2, Remark 9.2.6 and Theorem 9.3.21]
{fghLiu}. Their restrictions to the good elliptic
model have the distinct centres (A.60.9), while
the main valuation lies above its generic point.
Their genera already exhaust the generic genus:
\[
                              1+7\cdot21=148.          \tag{A.60.14}
\]
The genus formula forces every further component
to be rational and the dual graph to be a tree.
All component Jacobians have seven-rank zero,
including the supersingular elliptic main, so
the good special Jacobian has seven-rank zero.

The actual degree-seven map \(q\) extends to the
good elliptic model by the same resolution-and-descent
argument as A.58. It has degree seven on the
radicial main component. Constancy of the degree
of a pulled-back line bundle makes every other
component map constantly. Every off-main subtree
has one attachment point and maps to its image.

All forty-nine points of the generic \(E[7]\)
specialize to \(O\), since the elliptic special
fibre is supersingular. Their ramified \(q\)-fibres
are the forty-nine distinct cusps. None can
specialize on one of the seven positive-genus
tails or its connecting subtree, since all
those subtrees map to the finite points (A.60.9).
The radicial main has a unique point over \(O\).
After separating the cusp sections, all forty-nine
marks therefore lie on the rational marked tree
attached at that point. No mark can remain at
the main attachment node. Thus every selected
cusp has a rational joining path to another,
with rational endpoints.

The all-prime full fixed-inertia conclusion is
exactly A.41--A.42. The actual cover has a
seven-group geometric Galois closure, so its
all-prime normalizer conclusion follows from
A.59. Alternatively away from seven its proper
curve has potential good reduction by the tame
argument and A.57 applies, while at seven the
explicit zero-rank Jacobian gives A.46. No
characteristic-zero Jacobian decomposition or
unverified arithmetic character is required. The exact
certificate verifies (A.60.3)--(A.60.4), (A.60.7)--(A.60.13),
the missing degree-seven coefficient, all scaling valuations,
and the genus and cusp counts; the geometric conclusions
use the full model argument above.
\end{proof}
